\chapter{2023年北京卷}

\section{单选题}

\begin{problem}
已知集合 \(M=\{x\mid x+2\ge 0\}\)，\(N=\{x\mid x-1<0\}\)，则 \(M\cap N=\)（\quad）

\choices
    {\(\{x\mid -2\le x<1\}\)}
    {\(\{x\mid -2<x\le 1\}\)}
    {\(\{x\mid x\ge -2\}\)}
    {\(\{x\mid x<1\}\)}

\end{problem}

\begin{answer}
A
\end{answer}

\begin{solution}
由题意
\(M=\{x\mid x\ge -2\}\)，\(N=\{x\mid x<1\}\).
因此
\(M\cap N=\{x\mid -2\le x<1\}\).
故选 A.
\end{solution}

\begin{problem}
在复平面内，复数 \(z\) 对应的点的坐标是 \((-1,\sqrt{3})\)，则 \(z\) 的共轭复数 \(\overline{z}=\)（\quad）

\choices
    {\(1+\sqrt{3}\i\)}
    {\(1-\sqrt{3}\i\)}
    {\(-1+\sqrt{3}\i\)}
    {\(-1-\sqrt{3}\i\)}

\end{problem}

\begin{answer}
D
\end{answer}

\begin{solution}
由复数的几何意义，
\(z=-1+\sqrt{3}\i\).
因此
\(\overline{z}=-1-\sqrt{3}\i\).
故选 D.
\end{solution}

\begin{problem}
已知向量 \(\bs{a}\)，\(\bs{b}\) 满足 \(\bs{a}+\bs{b}=(2,3)\)，\(\bs{a}-\bs{b}=(-2,1)\)，则 \(|\bs{a}|^2-|\bs{b}|^2=\)（\quad）

\choices
    {\(-2\)}
    {\(-1\)}
    {0}
    {1}

\end{problem}

\begin{answer}
B
\end{answer}

\begin{solution}
由数量积公式，
\(|\bs{a}|^2-|\bs{b}|^2=(\bs{a}+\bs{b})\cdot(\bs{a}-\bs{b})\).
故
\(|\bs{a}|^2-|\bs{b}|^2=2\times(-2)+3\times 1=-1\).
故选 B.
\end{solution}

\begin{problem}
下列函数中，在区间 \((0,+\infty)\) 上单调递增的是（\quad）

\choices
    {\(f(x)=-\ln x\)}
    {\(f(x)=\dfrac{1}{2^x}\)}
    {\(f(x)=-\dfrac{1}{x}\)}
    {\(f(x)=3^{|x-1|}\)}

\end{problem}

\begin{answer}
C
\end{answer}

\begin{solution}
\(-\ln x\) 在 \((0,+\infty)\) 上单调递减，\(\dfrac{1}{2^x}=2^{-x}\) 在 \((0,+\infty)\) 上单调递减.
又
\(f(x)=-\frac{1}{x}\)
在 \((0,+\infty)\) 上单调递增.
而 \(3^{|x-1|}\) 在 \(x=1\) 两侧先减后增，不单调.
故选 C.
\end{solution}

\begin{problem}
\(\left(2x-\dfrac{1}{x}\right)^5\) 的展开式中 \(x\) 的系数为（\quad）

\choices
    {\(-80\)}
    {\(-40\)}
    {40}
    {80}

\end{problem}

\begin{answer}
D
\end{answer}

\begin{solution}
通项为
\(T_{r+1}=\mathrm{C}_5^r(2x)^{5-r}\left(-\frac{1}{x}\right)^r =(-1)^r 2^{5-r}\mathrm{C}_5^r x^{5-2r}\).
令 \(5-2r=1\)，得 \(r=2\). 故所求系数为
\((-1)^2 2^3\mathrm{C}_5^2=8\times 10=80\).
故选 D.
\end{solution}

\begin{problem}
已知抛物线 \(C:y^2=8x\) 的焦点为 \(F\)，点 \(M\) 在 \(C\) 上. 若 \(M\) 到直线 \(x=-3\) 的距离为 \(5\)，则 \(|MF|=\)（\quad）

\choices
    {7}
    {6}
    {5}
    {4}

\end{problem}

\begin{answer}
D
\end{answer}

\begin{solution}
抛物线 \(y^2=8x\) 的焦点为 \(F(2,0)\)，准线为 \(x=-2\).
由抛物线定义，点 \(M\) 到准线 \(x=-2\) 的距离等于 \(|MF|\). 又点 \(M\) 到直线 \(x=-3\) 的距离为 5，所以
\(|MF|+1=5\).
故
\(|MF|=4\).
故选 D.
\end{solution}

\begin{problem}
在 \(\triangle ABC\) 中，\((a+c)(\sin A-\sin C)=b(\sin A-\sin B)\)，则 \(\angle C=\)（\quad）

\choices
    {\(\dfrac{\pi}{6}\)}
    {\(\dfrac{\pi}{3}\)}
    {\(\dfrac{2\pi}{3}\)}
    {\(\dfrac{5\pi}{6}\)}

\end{problem}

\begin{answer}
B
\end{answer}

\begin{solution}
由正弦定理 \(\dfrac{a}{\sin A}=\dfrac{b}{\sin B}=\dfrac{c}{\sin C}\)，可将题设化为
\((a+c)(a-c)=b(a-b)\)，
即
\(a^2-c^2=ab-b^2\).
从而
\(a^2+b^2-c^2=ab\).
由余弦定理，
\(\cos C=\frac{a^2+b^2-c^2}{2ab}=\frac{1}{2}\).
又 \(0<C<\pi\)，故
\(C=\frac{\pi}{3}\).
故选 B.
\end{solution}

\begin{problem}
若 \(xy\ne 0\)，则“\(x+y=0\)”是“\(\dfrac{x}{y}+\dfrac{y}{x}=-2\)”的（\quad）

\choices
    {充分不必要条件}
    {必要不充分条件}
    {充要条件}
    {既不充分也不必要条件}

\end{problem}

\begin{answer}
C
\end{answer}

\begin{solution}
若
\(\frac{x}{y}+\frac{y}{x}=-2\)，
则
\(\frac{x^2+y^2}{xy}=-2\)，
即
\(x^2+2xy+y^2=0\)，
故
\((x+y)^2=0\)，
从而 \(x+y=0\).

反过来，若 \(x+y=0\)，则 \(y=-x\)，于是
\(\frac{x}{y}+\frac{y}{x}=\frac{x}{-x}+\frac{-x}{x}=-2\).
所以“\(x+y=0\)”是“\(\dfrac{x}{y}+\dfrac{y}{x}=-2\)”的充要条件.
故选 C.
\end{solution}

\begin{problem}
坡屋顶是我国传统建筑造型之一，蕴含着丰富的数学元素. 安装灯带可以勾勒出建筑轮廓，展现造型之美. 如图，某坡屋顶可视为一个五面体，其中两个面是全等的等腰梯形，两个面是全等的等腰三角形. 若 \(AB=25\text{m}\)，\(BC=AD=10\text{m}\)，且等腰梯形所在的平面，等腰三角形所在的平面与平面 \(ABCD\) 的夹角的正切值均为 \(\dfrac{\sqrt{14}}{5}\)，则该五面体的所有棱长之和为（\quad）

\bitmapfigure[max width=6.4cm]{img/2023/beijing/q9_fig1.png}

\choices
    {\(102\text{m}\)}
    {\(112\text{m}\)}
    {\(117\text{m}\)}
    {\(125\text{m}\)}

\end{problem}

\begin{answer}
C
\end{answer}

\begin{solution}
如图，过 \(E\) 作 \(EO\perp \text{平面 }ABCD\)，垂足为 \(O\)，再作 \(EG\perp BC\)，\(EM\perp AB\)，垂足分别为 \(G\)，\(M\).

\solutionfigure{\bitmapfigure[max width=6.6cm]{img/2023/beijing/q9_solution_fig1.png}}

由题意，等腰梯形所在平面与底面的夹角为 \(\angle EMO\)，等腰三角形所在平面与底面的夹角为 \(\angle EGO\)，故
\(\tan \angle EMO=\tan \angle EGO=\frac{\sqrt{14}}{5}\).
又由对称性可知
\(OM=OG=\frac{BC}{2}=5\)，
于是
\(EO=5\times \frac{\sqrt{14}}{5}=\sqrt{14}\).
在直角三角形 \(EOG\) 中，
\(EG=\sqrt{EO^2+OG^2}=\sqrt{14+25}=\sqrt{39}\).
在直角三角形 \(EBG\) 中，
\(EB=\sqrt{EG^2+BG^2}=\sqrt{39+25}=8\).
又 \(EF=AB-OM-OG=25-5-5=15\).
故所有棱长之和为
\(2AB+2BC+EF+4EB =2\times 25+2\times 10+15+4\times 8 =117\text{m}\).
故选 C.
\end{solution}

\begin{problem}
已知数列 \(\{a_n\}\) 满足 \(a_{n+1}=\frac{1}{4}(a_n-6)^3+6\)（\(n=1\)，\(2\)，\(3\)，\(\dots\)），则（\quad）

\choices
    {当 \(a_1=3\) 时，\(\{a_n\}\) 为递减数列，且存在常数 \(M\le 0\)，使得 \(a_n>M\) 恒成立}
    {当 \(a_1=5\) 时，\(\{a_n\}\) 为递增数列，且存在常数 \(M\le 6\)，使得 \(a_n<M\) 恒成立}
    {当 \(a_1=7\) 时，\(\{a_n\}\) 为递减数列，且存在常数 \(M>6\)，使得 \(a_n>M\) 恒成立}
    {当 \(a_1=9\) 时，\(\{a_n\}\) 为递增数列，且存在常数 \(M>0\)，使得 \(a_n<M\) 恒成立}

\end{problem}

\begin{answer}
B
\end{answer}

\begin{solution}
记
\(F(x)=\frac{1}{4}(x-6)^3+6-x =\frac{1}{4}(x-4)(x-6)(x-8)\).
则
\(a_{n+1}-a_n=F(a_n)\).
由上式可知：
\(F(x)<0\)，\(\text{当 }x<4\text{\ 或\ }6<x<8;\)，\(F(x)>0\)，\(\text{当 }4<x<6\text{\ 或\ }x>8\).

对 A，若 \(a_1=3\)，可由递推式归纳得 \(a_n\le 3<4\)，故 \(\{a_n\}\) 递减. 但对
\(H(x)=\frac{1}{4}x^3-\frac{9}{2}x^2+26x-47\)
有 \(H(x)<0\)（当 \(x\le 3\)），所以
\(a_{n+1}-a_n+1=H(a_n)<0\)，
即 \(a_{n+1}<a_n-1\). 于是数列无下界形如 \(M\le 0\) 使 \(a_n>M\) 恒成立，A 错.

对 B，若 \(a_1=5\)，由递推式易归纳得
\(5\le a_n<6\).
又 \(F(x)>0\) 在 \((4,6)\) 上成立，所以 \(\{a_n\}\) 递增；并且始终有 \(a_n<6\)，取 \(M=6\) 即可，B 对.

对 C，若 \(a_1=7\)，由递推式可归纳得
\(a_n=6+\left(\frac{1}{4}\right)^{\frac{3^{\,n-1}-1}{2}}\)（\(n\ge 2\)）.
故 \(6<a_n\le 7\)，且数列递减并趋于 6，所以不存在 \(M>6\) 使 \(a_n>M\) 恒成立，C 错.

对 D，若 \(a_1=9\)，归纳可得 \(a_n\ge 9\). 又 \(F(x)>0\) 在 \((8,+\infty)\) 上成立，故数列递增. 再由
\(a_{n+1}-a_n-1=\frac{1}{4}a_n^3-\frac{9}{2}a_n^2+26a_n-49>0\)（\(a_n\ge 9\)）
得 \(a_{n+1}>a_n+1\)，所以不可能存在正数 \(M\) 使 \(a_n<M\) 恒成立，D 错.
故选 B.
\end{solution}
\section{填空题}

\begin{problem}
已知函数 \(f(x)=4^x+\log_2 x\)，则 \(f\left(\dfrac{1}{2}\right)=\fillinblank{}\).

\end{problem}

\begin{answer}
1
\end{answer}

\begin{solution}
直接代入可得
\(f\left(\frac{1}{2}\right)=4^{\frac{1}{2}}+\log_2 \frac{1}{2}=2-1=1\).
\end{solution}

\begin{problem}
已知双曲线 \(C\) 的焦点为 \((-2,0)\) 和 \((2,0)\)，离心率为 \(\sqrt{2}\)，则 \(C\) 的方程为\fillinblank{}.

\end{problem}

\begin{answer}
\(\dfrac{x^2}{2}-\dfrac{y^2}{2}=1\)
\end{answer}

\begin{solution}
设双曲线方程为
\(\frac{x^2}{a^2}-\frac{y^2}{b^2}=1\).
则半焦距 \(c=2\)，又离心率
\(e=\frac{c}{a}=\sqrt{2}\)，
故
\(a=\sqrt{2}\).
由
\(c^2=a^2+b^2\)
得
\(b^2=4-2=2\).
所以双曲线方程为
\(\frac{x^2}{2}-\frac{y^2}{2}=1\).
\end{solution}

\begin{problem}
已知命题 \(p\)：若 \(\alpha\)，\(\beta\) 为第一象限角，且 \(\alpha>\beta\)，则 \(\tan\alpha>\tan\beta\). 能说明 \(p\) 为假命题的一组 \(\alpha\)，\(\beta\) 的值为 \(\alpha=\fillinblank{}\)，\(\beta=\fillinblank{}\).

\end{problem}

\begin{answer}
\(\dfrac{9\pi}{4}\)；\(\dfrac{\pi}{3}\)
\end{answer}

\begin{solution}
因为 \(\tan x\) 在 \(\left(0,\frac{\pi}{2}\right)\) 上单调递增，可取
\(\alpha_0=\frac{\pi}{4}\)，\(\beta_0=\frac{\pi}{3}\)，
则
\(0<\alpha_0<\beta_0<\frac{\pi}{2}\)，\(\tan\alpha_0<\tan\beta_0\).
再取
\(\alpha=2\pi+\alpha_0=\frac{9\pi}{4}\)，\(\beta=\beta_0=\frac{\pi}{3}\)，
则 \(\alpha>\beta\)，但
\(\tan\alpha=\tan\frac{\pi}{4}=1<\sqrt{3}=\tan\frac{\pi}{3}\).
故可填
\(\alpha=\frac{9\pi}{4}\)，\(\beta=\frac{\pi}{3}\).
\end{solution}

\begin{problem}
我国度量衡的发展有着悠久的历史，战国时期就已经出现了类似于砝码的，用来测量物体质量的“环权”. 已知 9 枚环权的质量（单位：铢）从小到大构成项数为 9 的数列 \(\{a_n\}\)，该数列的前 3 项成等差数列，后 7 项成等比数列，且 \(a_1=1\)，\(a_5=12\)，\(a_9=192\)，则 \(a_7=\fillinblank{}\)；数列 \(\{a_n\}\) 所有项的和为 \(\fillinblank{}\).

\end{problem}

\begin{answer}
48；384
\end{answer}

\begin{solution}
因为后 7 项成等比数列，所以 \(a_7^2=a_5a_9=12\times 192=48^2\).
又各项为正，故
\(a_7=48\).
再由
\(a_5^2=a_3a_7\)
得
\(a_3=\frac{a_5^2}{a_7}=\frac{144}{48}=3\).
设后 7 项公比为 \(q>0\)，则
\(q^2=\frac{a_5}{a_3}=\frac{12}{3}=4\)，
所以 \(q=2\).
前 3 项成等差数列，故 \(a_1+a_2+a_3=\frac{3(a_1+a_3)}{2}=\frac{3(1+3)}{2}=6\).
后 7 项和为
\[
a_3+a_4+\cdots+a_9 =\frac{a_3-a_9q}{1-q} =\frac{3-192\times 2}{1-2} =381
\]
因此总和为 \(6+381-a_3=6+381-3=384\).
故答案为 48，384.
\end{solution}

\begin{problem}
设 \(a>0\)，函数 \(f(x)=x+2\)（\(x<-a\)），\(f(x)=\sqrt{a^2-x^2}\)（\(-a\le x\le a\)），\(f(x)=-\sqrt{x}-1\)（\(x>a\)）. 给出下列四个结论：

\begin{circlelist}
\item \(f(x)\) 在区间 \((a-1,+\infty)\) 上单调递减；

\item 当 \(a\ge 1\) 时，\(f(x)\) 存在最大值；

\item 设 \(M(x_1,f(x_1))(x_1\le a)\)，\(N(x_2,f(x_2))(x_2>a)\)，则 \(|MN|>1\)；

\item 设 \(P(x_3,f(x_3))(x_3<-a)\)，\(Q(x_4,f(x_4))(x_4\ge -a)\).
\end{circlelist}
若 \(|PQ|\) 存在最小值，则 \(a\) 的取值范围是 \(\left(0,\dfrac{1}{2}\right]\).

其中所有正确结论的序号是\fillinblank{}.

\end{problem}

\begin{answer}
\(\circled{2}\circled{3}\)
\end{answer}

\begin{solution}
当 \(x<-a\) 时，图像是一条斜率为 1 的射线；当 \(-a\le x\le a\) 时，图像是半径为 \(a\) 的上半圆；当 \(x>a\) 时，图像是 \(y=-\sqrt{x}-1\) 的递减曲线.

对 \(\circled{1}\)，取 \(a=\dfrac{1}{2}\)，图像如下：

\solutionfigure{\bitmapfigure[max width=4.0cm]{img/2023/beijing/q15_solution_fig1.png}}

此时 \(a-1=-\dfrac{1}{2}\)，在区间 \(\left(-\dfrac{1}{2},0\right)\) 上，半圆部分单调递增，故 \(\circled{1}\) 错.

对 \(\circled{2}\)，当 \(a\ge 1\) 时，
\(x<-a \implies f(x)=x+2<-a+2\le 1\)，
\(-a\le x\le a \implies f(x)=\sqrt{a^2-x^2}\le a\)，
且在 \(x=0\) 处取到 \(a\)；
\(x>a \implies f(x)=-\sqrt{x}-1<-\sqrt{a}-1\le -2\).
因此 \(f(x)\) 存在最大值 \(a\)，故 \(\circled{2}\) 对.

对 \(\circled{3}\)，结合图像可知，当 \(x_1=a\) 且 \(x_2\) 从右侧逼近 \(a\) 时，\(|MN|\) 最小，如下图：

\solutionfigure{\bitmapfigure[max width=4.0cm]{img/2023/beijing/q15_solution_fig2.png}}

此时
\(y_1=f(a)=0\)，\(y_2<-\sqrt{a}-1\)，
故
\(|MN|>y_1-y_2>\sqrt{a}+1>1\).
所以 \(\circled{3}\) 对.

对 \(\circled{4}\)，取 \(a=\dfrac{4}{5}\)，图像如下：

\solutionfigure{\bitmapfigure[max width=4.3cm]{img/2023/beijing/q15_solution_fig3.png}}

由图可知，若要 \(|PQ|\) 取得最小值，则点 \(P\) 在线 \(y=x+2\) 上，点 \(Q\) 在半圆上，且连接线与该直线垂直. 直线 \(y=x+2\) 的垂线过原点可取为 \(y=-x\)，联立 \(y=-x\) 与 \(y=x+2\)，得 \(P(-1,1)\).
此点满足 \(-1<-\dfrac{4}{5}\)，故 \(a=\dfrac{4}{5}\) 时 \(|PQ|\) 也存在最小值，与“\(a\) 的范围仅为 \(\left(0,\dfrac{1}{2}\right]\)”矛盾，所以 \(\circled{4}\) 错.
故答案为 \(\circled{2}\circled{3}\).
\end{solution}
\section{解答题}

\begin{problem}
如图，在三棱锥 \(P-ABC\) 中，\(PA\perp \text{平面 }ABC\)，\(PA=AB=BC=1\)，\(PC=\sqrt{3}\).

\begin{enumerate}
\item 求证：\(BC\perp \text{平面 }PAB\)；
\item 求二面角 \(A-PC-B\) 的大小.
\end{enumerate}

\bitmapfigure[max width=4.9cm]{img/2023/beijing/q16_fig1.png}

\end{problem}

\begin{solution}
（1）因为 \(PA\perp \text{平面 }ABC\)，且 \(BC\subset \text{平面 }ABC\)，所以
\(PA\perp BC\).
又 \(AB=1\)，\(BC=1\)，\(PA=1\)，而
\(PB=\sqrt{PA^2+AB^2}=\sqrt{2}\)，
从而 \(PB^2+BC^2=2+1=3=PC^2\).
故 \(\triangle PBC\) 为直角三角形，且
\(BC\perp PB\).
于是 \(BC\) 同时垂直于平面 \(PAB\) 内两条相交直线 \(PA\)，\(PB\)，故
\(BC\perp \text{平面 }PAB\).

（2）由（1）知 \(BC\perp AB\). 以 \(A\) 为原点，\(AB\) 所在直线为 \(x\) 轴，过 \(A\) 且与 \(BC\) 平行的直线为 \(y\) 轴，\(AP\) 为 \(z\) 轴，建立空间直角坐标系.

\solutionfigure{\bitmapfigure[max width=4.9cm]{img/2023/beijing/q16_solution_fig1.png}}

则
\(A(0,0,0)\)，\(P(0,0,1)\)，\(B(1,0,0)\)，\(C(1,1,0)\).
因此 \(\overrightarrow{AP}=(0,0,1)\)，\(\overrightarrow{AC}=(1,1,0)\)，\(\overrightarrow{BC}=(0,1,0)\)，\(\overrightarrow{PC}=(1,1,-1)\).
设平面 \(PAC\) 的法向量为 \(\bs{m}=(x_1,y_1,z_1)\)，则由 \(\bs{m}\cdot \overrightarrow{AP}=0\) 和 \(\bs{m}\cdot \overrightarrow{AC}=0\)，可取 \(\bs{m}=(1,-1,0)\).
设平面 \(PBC\) 的法向量为 \(\bs{n}=(x_2,y_2,z_2)\)，则由 \(\bs{n}\cdot \overrightarrow{BC}=0\) 和 \(\bs{n}\cdot \overrightarrow{PC}=0\)，可取 \(\bs{n}=(1,0,1)\).
于是
\[
\cos \angle(\bs{m},\bs{n}) =\frac{\bs{m}\cdot \bs{n}}{|\bs{m}||\bs{n}|} =\frac{1}{\sqrt{2}\cdot \sqrt{2}} =\frac{1}{2}
\]
故二面角 \(A-PC-B\) 的大小为 \(\frac{\pi}{3}\).
\end{solution}

\begin{problem}
设函数 \(f(x)=\sin \omega x\cos \varphi+\cos \omega x\sin \varphi\)（\(\omega>0\)，\(|\varphi|<\frac{\pi}{2}\)）.

\begin{enumerate}
\item 若 \(f(0)=-\dfrac{\sqrt{3}}{2}\)，求 \(\varphi\) 的值；
\item 已知 \(f(x)\) 在区间 \(\left[-\dfrac{\pi}{3},\dfrac{2\pi}{3}\right]\) 上单调递增，\(f\left(\dfrac{2\pi}{3}\right)=1\)，再从条件 \(\circled{1}\)，条件 \(\circled{2}\)，条件 \(\circled{3}\) 这三个条件中选择一个作为已知，使函数 \(f(x)\) 存在，求 \(\omega\)，\(\varphi\) 的值.
\end{enumerate}

条件 \(\circled{1}\)：\(f\left(\dfrac{\pi}{3}\right)=\sqrt{2}\)；

条件 \(\circled{2}\)：\(f\left(-\dfrac{\pi}{3}\right)=-1\)；

条件 \(\circled{3}\)：\(f(x)\) 在区间 \(\left[-\dfrac{\pi}{2},-\dfrac{\pi}{3}\right]\) 上单调递减.

\end{problem}

\begin{solution}
（1）由题意
\(f(0)=\sin \varphi=-\frac{\sqrt{3}}{2}\).
又 \(|\varphi|<\dfrac{\pi}{2}\)，故
\(\varphi=-\frac{\pi}{3}\).

（2）由两角和公式，
\(f(x)=\sin(\omega x+\varphi)\).
因此 \(f(x)\) 的最大值为 1，最小值为 \(-1\).

若选条件 \(\circled{1}\)，则
\(f\left(\frac{\pi}{3}\right)=\sqrt{2}>1\)，
不可能成立，所以条件 \(\circled{1}\) 不可取.

若选条件 \(\circled{2}\)，因为 \(f(x)\) 在 \(\left[-\dfrac{\pi}{3},\dfrac{2\pi}{3}\right]\) 上单调递增，且
\(f\left(-\frac{\pi}{3}\right)=-1\)，\(f\left(\frac{2\pi}{3}\right)=1\)，
故这两个点恰对应一个半周期，于是
\(\frac{T}{2}=\frac{2\pi}{3}-\left(-\frac{\pi}{3}\right)=\pi\).
从而
\(T=2\pi\)，\(\omega=\frac{2\pi}{T}=1\).
此时
\(f(x)=\sin(x+\varphi)\).
又
\(f\left(-\frac{\pi}{3}\right)=-1\)，
故
\(\sin\left(-\frac{\pi}{3}+\varphi\right)=-1\).
结合 \(|\varphi|<\dfrac{\pi}{2}\)，得
\(\varphi=-\frac{\pi}{6}\).

若选条件 \(\circled{3}\)，由“在 \(\left[-\dfrac{\pi}{3},\dfrac{2\pi}{3}\right]\) 上单调递增”与“在 \(\left[-\dfrac{\pi}{2},-\dfrac{\pi}{3}\right]\) 上单调递减”可知，\(x=-\dfrac{\pi}{3}\) 处取得最小值，即
\(f\left(-\frac{\pi}{3}\right)=-1\).
于是化为条件 \(\circled{2}\) 的情形，同样得到
\(\omega=1\)，\(\varphi=-\frac{\pi}{6}\).
\end{solution}

\begin{problem}
为研究某种农产品价格变化的规律，收集得到了该农产品连续 40 天的价格变化数据，如下表所示. 在描述价格变化时，用“+”表示“上涨”，即当天价格比前一天价格高；用“-”表示“下跌”，即当天价格比前一天价格低；用“0”表示“不变”，即当天价格与前一天价格相同. 用频率估计概率.

\begin{center}
\resizebox{\linewidth}{!}{
\begin{tabular}{|c|*{20}{c|}}
\hline
第 1 天到第 20 天 & - & + & + & 0 & - & - & - & + & + & 0 & + & 0 & - & - & + & - & + & 0 & 0 & + \\
\hline
第 21 天到第 40 天 & 0 & + & + & 0 & - & - & - & + & + & 0 & + & 0 & + & - & - & - & + & 0 & - & + \\
\hline
\end{tabular}
}
\end{center}

\begin{enumerate}
\item 试估计该农产品价格“上涨”的概率；
\item 假设该农产品每天的价格变化是相互独立的. 在未来的日子里任取 4 天，试估计该农产品价格在这 4 天中 2 天“上涨”，1 天“下跌”，1 天“不变”的概率；
\item 假设该农产品每天的价格变化只受前一天价格变化的影响. 判断第 41 天该农产品价格“上涨”“下跌”和“不变”的概率估计值哪个最大. （结论不要求证明）
\end{enumerate}

\end{problem}

\begin{solution}
（1）在 40 天中，“+” 共出现 16 次，因此价格“上涨”的概率估计值为
\(\frac{16}{40}=0.4\).
（2）由表中数据可得：
\(P(\text{上涨})=0.4\)，\(P(\text{下跌})=0.35\)，\(P(\text{不变})=0.25\).
故所求概率为
\(\mathrm{C}_4^2\times 0.4^2\times \mathrm{C}_2^1\times 0.35\times 0.25=0.168\).
（3）第 40 天处于“上涨”状态. 从前 39 次的 15 次“上涨”进行分析，“上涨”后下一次仍“上涨”的有 4 次，“下跌”的有 2 次，“不变”的有 9 次.
因此三者中概率估计值最大的是“不变”.
\end{solution}

\begin{problem}
已知椭圆 \(E:\frac{x^2}{a^2}+\frac{y^2}{b^2}=1\)（\(a>b>0\)） 的离心率为 \(\dfrac{\sqrt{5}}{3}\)，\(A\)，\(C\) 分别是 \(E\) 的上，下顶点，\(B\)，\(D\) 分别是 \(E\) 的左，右顶点，\(|AC|=4\).

\begin{enumerate}
\item 求 \(E\) 的方程；
\item 设 \(P\) 为第一象限内 \(E\) 上的动点，直线 \(PD\) 与直线 \(BC\) 交于点 \(M\)，直线 \(PA\) 与直线 \(y=-2\) 交于点 \(N\). 求证：\(MN\parallel CD\).
\end{enumerate}

\end{problem}

\begin{solution}
（1）由 \(|AC|=4\) 知
\(2b=4\)，\(b=2\).
又离心率
\(e=\frac{c}{a}=\frac{\sqrt{5}}{3}\)，
故
\(c=\frac{\sqrt{5}}{3}a\).
由
\(c^2=a^2-b^2\)
得
\(\frac{5}{9}a^2=a^2-4\)，
解得
\(a^2=9\).
因此椭圆方程为
\(\frac{x^2}{9}+\frac{y^2}{4}=1\).

（2）由（1）知
\(A(0,2)\)，\(C(0,-2)\)，\(B(-3,0)\)，\(D(3,0)\).
设
\(P(m,n)\)（\(0<m<3\)，\(0<n<2\)），
则
\(\frac{m^2}{9}+\frac{n^2}{4}=1\).

\solutionfigure{\bitmapfigure[max width=6.0cm]{img/2023/beijing/q19_solution_fig1.png}}

由
\(k_{BC}=\frac{-2-0}{0-(-3)}=-\frac{2}{3}\)，
得直线 \(BC\) 方程为
\(y=-\frac{2}{3}x-2\).
又
\(k_{PD}=\frac{n-0}{m-3}=\frac{n}{m-3}\)，
故直线 \(PD\) 方程为
\(y=\frac{n}{m-3}(x-3)\).
联立可得
\(M\left(\frac{3(3n-2m+6)}{3n+2m-6},\ -\frac{12n}{3n+2m-6}\right)\).
再由
\(k_{PA}=\frac{n-2}{m}\)，
知直线 \(PA\) 方程为
\(y=\frac{n-2}{m}x+2\).
令 \(y=-2\)，得
\(N\left(-\frac{4m}{n-2},\ -2\right)\).
又
\(\frac{m^2}{9}+\frac{n^2}{4}=1\)，
则
\(m^2=9-\frac{9n^2}{4}\)，
从而
\(8m^2=72-18n^2\).
所以
\[
\begin{aligned}
k_{MN}
&=\frac{-\dfrac{12n}{3n+2m-6}+2}{\dfrac{3(3n-2m+6)}{3n+2m-6}+\dfrac{4m}{n-2}} \\
&=\frac{(-6n+4m-12)(n-2)}{(9n-6m+18)(n-2)+4m(3n+2m-6)} \\
&=\frac{-6n^2+4mn-8m+24}{9n^2+8m^2+6mn-12m-36} \\
&=\frac{-6n^2+4mn-8m+24}{9n^2+72-18n^2+6mn-12m-36} \\
&=\frac{-6n^2+4mn-8m+24}{-9n^2+6mn-12m+36} \\
&=\frac{2(-3n^2+2mn-4m+12)}{3(-3n^2+2mn-4m+12)}
=\frac{2}{3}.
\end{aligned}
\]
而
\(k_{CD}=\frac{0-(-2)}{3-0}=\frac{2}{3}\).
所以
\(k_{MN}=k_{CD}\).
因此 \(MN\) 与 \(CD\) 不重合，故
\(MN\parallel CD\).
\end{solution}

\begin{problem}
设函数 \(f(x)=x-x^3\e^{ax+b}\)，曲线 \(y=f(x)\) 在点 \((1,f(1))\) 处的切线方程为 \(y=-x+1\).

\begin{enumerate}
\item 求 \(a\)，\(b\) 的值；
\item 设函数 \(g(x)=f'(x)\)，求 \(g(x)\) 的单调区间；
\item 求 \(f(x)\) 的极值点个数.
\end{enumerate}

\end{problem}

\begin{solution}
（1）因为切线方程为 \(y=-x+1\)，所以
\(f(1)=0\)，\(f'(1)=-1\).
又
\(f(x)=x-x^3\e^{ax+b}\)，
故
\(f'(x)=1-(3x^2+ax^3)\e^{ax+b}\).
代入 \(x=1\) 得 \(1-\e^{a+b}=0\)，\(1-(3+a)\e^{a+b}=-1\).
解得
\(a=-1\)，\(b=1\).

（2）由（1）得
\(g(x)=f'(x)=1-(3x^2-x^3)\e^{-x+1}\).
于是
\(g'(x)=-x(x^2-6x+6)\e^{-x+1}\).
令
\(x^2-6x+6=0\)，
得
\(x_1=3-\sqrt{3}\)，\(x_2=3+\sqrt{3}\).
因为 \(\e^{-x+1}>0\)，故 \(g'(x)\) 的符号由 \(-x(x^2-6x+6)\) 决定，从而
\(g(x)\)
在
\((0,3-\sqrt{3})\)，\((3+\sqrt{3},+\infty)\)
上单调递减，在
\((-\infty,0)\)，\((3-\sqrt{3},3+\sqrt{3})\)
上单调递增.

（3）因为 \(f'(x)=g(x)\)，所以只需考察 \(g(x)=0\) 的零点个数.

在 \((-\infty,0)\) 上，\(g(x)\) 单调递增，且
\(g(-1)=1-4\e^2<0\)，\(g(0)=1>0\)，
故在 \((-\infty,0)\) 上有唯一零点.

在 \((0,3-\sqrt{3})\) 上，\(g(x)\) 单调递减，且
\(g(0)=1>0\)，\(g(1)=1-2<0\)，
故在 \((0,3-\sqrt{3})\) 上有唯一零点.

在 \((3-\sqrt{3},3+\sqrt{3})\) 上，\(g(x)\) 单调递增. 又由上一段可知 \(g(1)<0\)，且
\(1<3-\sqrt{3}\)，
故
\(g(3-\sqrt{3})<g(1)<0\).
又因为
\(3-\sqrt{3}<3<3+\sqrt{3}\)，\(g(3)=1>0\)，
所以 \(g(x)\) 在 \((3-\sqrt{3},3+\sqrt{3})\) 上有唯一零点.

而当 \(x>3+\sqrt{3}>3\) 时，
\(3x^2-x^3=x^2(3-x)<0\)，
于是
\(g(x)=1-(3x^2-x^3)\e^{-x+1}>0\)，
故在 \((3+\sqrt{3},+\infty)\) 上无零点.

综上，\(f'(x)\) 共有 3 个零点，所以 \(f(x)\) 有 3 个极值点.
\end{solution}

\begin{problem}
已知数列 \(\{a_n\}\)，\(\{b_n\}\) 的项数均为 \(m\ (m>2)\)，且 \(a_n\)，\(b_n\in\{1,2,\dots,m\}\)，\(\{a_n\}\)，\(\{b_n\}\) 的前 \(n\) 项和分别为 \(A_n\)，\(B_n\)，并规定 \(A_0=B_0=0\). 对于 \(k\in\{0,1,2,\dots,m\}\)，定义 \(r_k=\max \{ i\mid B_i\le A_k,\ i\in\{0,1,2,\dots,m\}\}\)，其中，\(\max M\) 表示数集 \(M\) 中最大的数.

\begin{enumerate}
\item 若 \(a_1=2\)，\(a_2=1\)，\(a_3=3\)，\(b_1=1\)，\(b_2=3\)，\(b_3=3\)，求 \(r_0\)，\(r_1\)，\(r_2\)，\(r_3\) 的值；
\item 若 \(a_1\ge b_1\)，且 \(2r_j\le r_{j+1}+r_{j-1}\ (j=1,2,\dots,m-1)\)，求 \(r_n\)；
\item 证明：存在 \(p\)，\(q\)，\(s\)，\(t\in\{0,1,2,\dots,m\}\)，满足 \(p>q\)，\(s>t\)，使得 \(A_p+B_t=A_q+B_s\).
\end{enumerate}

\end{problem}

\begin{solution}
（1）由题意，\(A_0=0\)，\(A_1=2\)，\(A_2=3\)，\(A_3=6\)，且 \(B_0=0\)，\(B_1=1\)，\(B_2=4\)，\(B_3=7\).
于是 \(r_0=0\)，\(r_1=1\)，\(r_2=1\)，\(r_3=2\).

（2）设
\(d_j=r_j-r_{j-1}\)（\(j=1\)，\(2\)，\(\dots\)，\(m\)）.
由条件
\(2r_j\le r_{j+1}+r_{j-1}\)
可得
\(d_{j+1}\ge d_j\).
即数列 \(\{d_j\}\) 单调不减.
又由于 \(a_1\ge b_1\)，故
\(A_1\ge B_1>B_0\)，
从而
\(r_1\ge 1\)，\(r_0=0\)，
所以
\(d_1=r_1-r_0\ge 1\).
若存在某个 \(d_j\ge 2\)，则由单调不减性知
\(d_j\)，\(d_{j+1}\)，\(\dots\)，\(d_m\ge 2\).
于是
\[
r_m=\sum_{k=1}^m d_k \ge (j-1)\times 1 + 2(m-j+1) =2m-j+1>m
\]
这与 \(r_m\le m\) 矛盾.
故所有 \(d_j=1\). 于是
\(r_n=r_0+\sum_{k=1}^n d_k=n\).

（3）分三种情况.

若 \(A_m=B_m\)，取
\(q=t=0\)，\(p=s=m\)，
则
\(A_p+B_t=A_m=A_q+B_s\).

若 \(A_m<B_m\)，对 \(1\le n\le m\) 定义
\(S_n=B_{r_n}-A_n\).
由 \(r_n\) 的定义知 \(S_n\le 0\)，且 \(S_n\) 为整数.
另一方面，若 \(S_n\le -m\)，则
\(B_{r_n+1}-A_n>0\)，
从而
\(b_{r_n+1}=B_{r_n+1}-B_{r_n} =(B_{r_n+1}-A_n)-(B_{r_n}-A_n)>m\)，
与 \(b_{r_n+1}\in\{1,2,\dots,m\}\) 矛盾.
故
\(1-m\le S_n\le 0\).

若存在 \(N\) 使 \(S_N=0\)，则
\(A_N=B_{r_N}\).
取
\(q=t=0\)，\(p=N\)，\(s=r_N\)，
则满足 \(p>q\)，\(s>t\)，\(A_p+B_t=A_q+B_s\).

若所有 \(S_n\ne 0\)，则
\(S_n\in\{-1,-2,\dots,-(m-1)\}\).
共有 \(m\) 个数 \(S_1\)，\(\dots\)，\(S_m\)，却只有 \(m-1\) 个可能取值，故存在 \(1\le X<Y\le m\) 使
\(S_X=S_Y\).
于是
\(B_{r_X}-A_X=B_{r_Y}-A_Y\)，
即
\(A_Y+B_{r_X}=A_X+B_{r_Y}\).
又因 \(A_Y>A_X\)，必有 \(B_{r_Y}>B_{r_X}\)，故 \(r_Y>r_X\).
取 \(p=Y\)，\(q=X\)，\(s=r_Y\)，\(t=r_X\)，
则满足 \(p>q\)，\(s>t\)，\(A_p+B_t=A_q+B_s\).

若 \(A_m>B_m\)，类似地定义
\(R_n=\max\{i\mid A_i\le B_n,\ i\in\{0,1,2,\dots,m\}\}\)，\(T_n=A_{R_n}-B_n\).
若存在 \(N\) 使 \(T_N=0\)，则 \(A_{R_N}=B_N\). 取
\(q=t=0\)，\(p=R_N\)，\(s=N\)，
即可得到 \(p>q\)，\(s>t\)，\(A_p+B_t=A_q+B_s\).
若所有 \(T_n\ne 0\)，则 \(T_n\in\{-1,-2,\dots,-(m-1)\}\). 共有 \(m\) 个数 \(T_1\)，\(\dots\)，\(T_m\)，却只有 \(m-1\) 个可能取值，故存在 \(1\le X<Y\le m\) 使 \(T_X=T_Y\). 于是
\(A_{R_X}-B_X=A_{R_Y}-B_Y\)，
即
\(A_{R_Y}+B_X=A_{R_X}+B_Y\).
又因 \(A_{R_Y}>A_{R_X}\)，必有 \(B_Y>B_X\)，故 \(R_Y>R_X\).
取 \(p=R_Y\)，\(q=R_X\)，\(s=Y\)，\(t=X\)，
则满足 \(p>q\)，\(s>t\)，\(A_p+B_t=A_q+B_s\).

综上，命题成立.
\end{solution}
